% !PRINTED 379
\weberpage{379}
%<<CONT
\weq{14}{\begin{aligned}
a &= a_1\delta^2 + b_1\beta\delta - c_1\beta^2\\
b &= 2a_1\gamma\delta + b_1(\alpha\delta + \beta\gamma) - 2c_1\alpha\beta,\\
-c &= a_1\gamma^2 + b_1\alpha\gamma - c_1\alpha^2,
\end{aligned}}
und ferner
\weq{15}{b^2 + 4ac = b_1^2 + 4a_1c_1 = D.}
Da $a$, $b$, $c$ ohne gemeinsamen Theiler sind, so haben auch $a_1$, $b_1$, $c_1$ keinen gemeinsamen Theiler, wie man aus (14) ersieht, und $a_1$, $b_1$, $c_1$ haben dieselbe Bedeutung f\"ur $\omega_1$, wie $a$, $b$, $c$ f\"ur $\omega$.

\gesperrt{VII. Man sieht ferner aus (15), dass bei \"aquivalenten Zahlen nicht nur die Irrationalit\"at $\sqrt{d}$, sondern auch die Discriminante dieselbe ist.}

Setzt man
\weq{16}{\omega_1 = x_1 + y_1\sqrt{d} = \frac{\alpha(x+y\sqrt{d})+\beta}{\gamma(x+y\sqrt{d})+\delta},}
und erweitert, um $x_1$, $y_1$ zu bestimmen, den letzten Bruch mit $\gamma(x-y\sqrt{d})+\delta$, so ergiebt sich nach (4) und (13)
\weq{17}{2x_1 = \frac{b_1}{c_1}, \qquad y_1 = \frac{\varepsilon y c}{c_1},}
also wenn die $\sqrt{D}$ ebenso verstanden wird wie in (7) und (8), nach (9):
\weq{18}{\omega_1 = \frac{b_1 + \varepsilon\sqrt{D}}{2c_1}.}

\begin{center}
\S.\ 123.

\gesperrt{Reducirte Zahlen mit negativer Discriminante.} % [sic] section number \S.\,123 repeats the number already used earlier on this same page for \textquotedbl{}Reducirte Zahlen\textquotedbl{} as printed
\end{center}

Die bisherigen Betrachtungen sind gleichm\"assig auf die beiden F\"alle der reellen und imagin\"aren quadratischen Irrationalzahlen, also auf positive und negative Discriminanten anwendbar. Jetzt aber m\"ussen beide F\"alle von einander getrennt werden, und wir beginnen mit den negativen Discriminanten. Es handelt sich also um imagin\"are Zahlen:
\weq{1}{\omega = x + y\sqrt{d} = \xi + i\eta,}
worin $\xi$ und $\eta$ reell sind; es gen\"ugt, von den beiden conjugirten Zahlen $\xi+i\eta$, $\xi-i\eta$ die eine zu betrachten. Wir nehmen also $\eta$ positiv an und setzen folgende Definition fest:
%CONT>>
